\section{Largest BST in a Binary Tree}
\label{sec:bst-largest-bst}

\problemstatement{Given the root of a (general, not necessarily search)
binary tree, find the size (number of nodes) of the \textbf{largest
subtree} that is a valid BST.}

\begin{patternbox}
This closing problem of the chapter is a deliberate capstone: it combines
Validate BST's (Section~\ref{sec:bst-validate}) range-checking idea with
a \textbf{bottom-up, tree-DP} style of recursion -- each subtree must
report enough information upward (is it a valid BST? how big? what are
its min and max values?) for its \emph{parent} to decide, in $O(1)$
additional work, whether combining it with its sibling and itself forms
an even larger valid BST. The keyword combination \emph{``largest
substructure satisfying a property''} within a \emph{general} tree (not
already known to be a BST) is the signal for this bottom-up,
information-passing recursion style, much like dynamic programming on
trees.
\end{patternbox}

\subsection*{Understanding the problem carefully}

Section~\ref{sec:bst-validate} validated a BST top-down, passing a valid
range \emph{downward}. Here, we do not know in advance which subtrees are
valid BSTs, so we instead compute, for \emph{every} subtree, bottom-up:
whether it is itself a valid BST, its size if so, and its minimum and
maximum values (needed by its parent to check the BST invariant one
level up). A node's own subtree is a valid BST exactly when both children
report ``yes, I am a valid BST,'' \emph{and} the node's value is strictly
greater than its left child's reported maximum, \emph{and} strictly less
than its right child's reported minimum.

\begin{ideabox}[Bottom-Up: Each Subtree Reports (isBST, size, min, max)]
Recurse postorder (process both children before combining at the current
node). For a null node, return $(\textit{isBST}=\text{true},
\textit{size}=0, \textit{min}=+\infty, \textit{max}=-\infty)$ (a
neutral, vacuously-true base case that combines safely with any actual
node). At a real node: if both children report $\textit{isBST}=\text{true}$,
and $\textit{node.val} > \textit{leftInfo.max}$, and
$\textit{node.val} < \textit{rightInfo.min}$: this subtree is a valid
BST of size $\textit{leftInfo.size} + \textit{rightInfo.size} + 1$; update
a running global answer with this size, and return
$(\text{true}, \textit{size}, \min(\textit{leftInfo.min},
\textit{node.val}), \max(\textit{rightInfo.max}, \textit{node.val}))$.
Otherwise, this subtree is not a valid BST: return
$\textit{isBST}=\text{false}$ (its size/min/max no longer matter to any
ancestor, since no ancestor can validly combine across an already-invalid
subtree).
\end{ideabox}

\subsection*{Why the global answer must be tracked throughout, not just at the root}

The largest valid BST is not necessarily the \emph{whole} tree, nor even
necessarily a subtree rooted near the root -- it could be a valid BST
entirely contained within one branch, while everything above that branch
fails the BST invariant. Because the recursion already computes,
\emph{at every single node}, whether that node's subtree is a valid BST
and how large it is, the correct overall answer is simply the
\textbf{maximum} of this size over every node visited during the entire
traversal -- which is why a running global maximum, updated every time a
valid BST subtree is found (regardless of whether its ancestors turn out
to be valid), is necessary, rather than only inspecting the final result
returned at the root.

\subsection*{Algorithm}

\begin{enumerate}
  \item Initialize a global $\textit{maxBSTSize} \gets 0$.
  \item $\textsc{Solve}(\textit{node})$:
  \begin{enumerate}
    \item If $\textit{node} = $ null: return $(\text{true}, 0, +\infty,
          -\infty)$.
    \item $\textit{L} \gets \textsc{Solve}(\textit{node}.\textit{left})$,
          $\textit{R} \gets
          \textsc{Solve}(\textit{node}.\textit{right})$.
    \item If $\textit{L}.\textit{isBST}$ and $\textit{R}.\textit{isBST}$
          and $\textit{node}.\textit{val} > \textit{L}.\textit{max}$ and
          $\textit{node}.\textit{val} < \textit{R}.\textit{min}$:
    \begin{itemize}
      \item $\textit{size} \gets \textit{L}.\textit{size} +
            \textit{R}.\textit{size} + 1$.
      \item $\textit{maxBSTSize} \gets \max(\textit{maxBSTSize},
            \textit{size})$.
      \item Return $(\text{true}, \textit{size},
            \min(\textit{L}.\textit{min}, \textit{node}.\textit{val}),
            \max(\textit{R}.\textit{max}, \textit{node}.\textit{val}))$.
    \end{itemize}
    \item Else: return $(\text{false}, 0, -, -)$.
  \end{enumerate}
  \item Call $\textsc{Solve}(\textit{root})$; return
        $\textit{maxBSTSize}$.
\end{enumerate}

\subsection*{Dry run}

\dryrun{Take the tree with root $5$, left child $2$ (left child $1$,
right child $3$, a valid BST of size $3$), right child $7$ (left child
$6$, right child $4$ -- note $4$ on the right side violates BST order).}

\begin{itemize}
  \item $\textsc{Solve}(1)$: leaf. Returns $(\text{true},1,1,1)$.
  \item $\textsc{Solve}(3)$: leaf. Returns $(\text{true},1,3,3)$.
  \item $\textsc{Solve}(2)$: $L=(\text{true},1,1,1)$,
        $R=(\text{true},1,3,3)$. Check: $2>1$ and $2<3$: valid. Size
        $=1+1+1=3$. $\textit{maxBSTSize}=3$. Return
        $(\text{true},3,1,3)$.
  \item $\textsc{Solve}(6)$: leaf. Returns $(\text{true},1,6,6)$.
  \item $\textsc{Solve}(4)$: leaf. Returns $(\text{true},1,4,4)$.
  \item $\textsc{Solve}(7)$: $L=(\text{true},1,6,6)$ (left child $6$),
        $R=(\text{true},1,4,4)$ (right child $4$). Check: $7>6$? Yes.
        $7<4$? No -- invalid (right child's minimum $4$ is not greater
        than $7$). Return $(\text{false}, -, -, -)$.
  \item $\textsc{Solve}(5)$: $L=(\text{true},3,1,3)$ (from node $2$),
        $R=(\text{false}, \ldots)$ (from node $7$). $R.\textit{isBST}$
        is false, so node $5$'s subtree is not a valid BST. Return
        $(\text{false}, -, -, -)$.
\end{itemize}

The largest valid BST found during the traversal has size $3$ (the
subtree rooted at node $2$), even though the overall tree (rooted at $5$)
is not a valid BST.

\subsection*{C++ implementation}

\begin{lstlisting}
#include <bits/stdc++.h>
using namespace std;

struct TreeNode {
    int val;
    TreeNode* left;
    TreeNode* right;
    TreeNode(int v) : val(v), left(nullptr), right(nullptr) {}
};

struct Info {
    bool isBST;
    int size;
    int minVal, maxVal;
};

int maxBSTSize = 0;

Info solve(TreeNode* node) {
    if (!node) return {true, 0, INT_MAX, INT_MIN};

    Info L = solve(node->left);
    Info R = solve(node->right);

    if (L.isBST && R.isBST && node->val > L.maxVal && node->val < R.minVal) {
        int size = L.size + R.size + 1;
        maxBSTSize = max(maxBSTSize, size);
        return {true, size, min(L.minVal, node->val), max(R.maxVal, node->val)};
    }
    return {false, 0, 0, 0};
}

int largestBSTSubtree(TreeNode* root) {
    maxBSTSize = 0;
    solve(root);
    return maxBSTSize;
}

int main() {
    TreeNode* root = new TreeNode(5);
    root->left = new TreeNode(2);
    root->left->left = new TreeNode(1);
    root->left->right = new TreeNode(3);
    root->right = new TreeNode(7);
    root->right->left = new TreeNode(6);
    root->right->right = new TreeNode(4);

    cout << "Largest BST size: " << largestBSTSubtree(root) << endl;
    return 0;
}
\end{lstlisting}

\begin{complexitybox}
\textbf{Time Complexity.} Each node is visited exactly once, doing
$O(1)$ work combining its children's reported info, giving
\[
  O(n).
\]
This is a substantial improvement over the naive approach of running
Validate BST (Section~\ref{sec:bst-validate}) independently on every
single subtree, which would cost $O(n^2)$.
\textbf{Space Complexity.} $O(h)$ for the recursion stack.
\end{complexitybox}

\begin{takeawaybox}
Whenever a problem asks for the largest/best substructure satisfying some
recursively-checkable property within a general tree, look for a
bottom-up recursion where each subtree reports exactly the information
its parent needs to make an $O(1)$ decision -- validity, size, and
boundary values, in this case. This ``return a small summary struct
upward, combine in $O(1)$ at each parent'' technique is the tree
analogue of dynamic programming, and it is the most general and reusable
idea in this entire chapter.
\end{takeawaybox}
